\section{Spreading metric for $\rho$-separators with terminals}\label{app:terminal-extension}

This appendix gives the extension of the procedure and analysis of Even et al.~\cite{FastApproximateGraphPartitioningAlgorithms} which gives an $O\br{\log n}$-approximation algorithm for the \textsc{$\rho$-SEP} problem, from the case $X=V$ to a designated terminal set $X\subseteq V$. The algorithm is almost exactly the same, and is based on a region-growing procedure. Whereas the original procedure grows a region around an arbitrary vertex, our modification procedure grows a region around an arbitrary terminal. The rest of the analysis is largely the same. 

\subsection{LP formulation}

For edge lengths $x$ and a vertex set $U\subseteq V$, we write $\dist_{x,U}(a,b)$ for the shortest-path distance between $a$ and $b$ in the induced subgraph $G[U]$ with edge lengths given by $x$.

Recall, the following is the weighted spreading-metrics \textsc{LP} relaxation \textsc{LP-$\rho$-SEP} for the \textsc{$\rho$-SEP} problem with a terminal set $X$:
\begin{align}
    \min \quad &\sum_{e\in E} c(e)\cdot x_e \\
    \text{subject to} \quad &\sum_{u\in S} \dist_{x,V}(v,u)\cdot w(u) \geq w\br{S}-\rho\cdot w\br{G}, \quad \forall S\subseteq V,\; v\in S\cap X, \label{eq:px1-main}\\
    &0\le x_e\le 1, \quad \forall e\in E.
\end{align}

Again, the LP can be solved in polynomial time using the ellipsoid method, since there is a polynomial-time separation oracle for the constraints \eqref{eq:px1-main}. Let $\tau_X$ denote the optimal value of \textsc{LP-$\rho$-SEP}.

\subsection{Lower bound}
Firstly, we claim that $\tau_X$ is a lower bound on the minimum capacity of a $(\rho,X)$-separator.
\begin{lemma}
$\tau_X$ is a lower bound on the minimum capacity of a $(\rho,X)$-separator.
\end{lemma}

\begin{proof}
Let $F$ be any $(\rho,X)$-separator. Define an integral length assignment by
\[
    d(e)=\begin{cases}
        1,& e\in F,\\
        0,& e\notin F.
    \end{cases}
\]
Take any constraint \eqref{eq:px1-main}, i.e., any $S\subseteq V$ and any $v\in S\cap X$.
Let $C$ be the connected component of $v$ in $G'=(V,E\setminus F)$.
For every $u\in S\setminus C$, every $v$-$u$ path uses at least one edge of $F$, hence $\dist_{x,V}(v,u)\ge 1$.
Therefore,
\begin{align*}
\sum_{u\in S}\dist_{x,V}(v,u)\cdot w(u)
&\ge \sum_{u\in S\setminus C}\dist_{x,V}(v,u)\cdot w(u)\\
&\ge w\br{S\setminus C}\\
&= w\br{S}-w\br{C\cap S}\\
&\ge w\br{S}-w\br{C}.
\end{align*}
Since $v\in X$, the component $C$ contains a terminal, so by definition of $(\rho,X)$-separator,
$w(C)\le \rho\cdot w\br{G}$. Hence
\[
\sum_{u\in S}\dist_{x,V}(v,u)\cdot w(u)\ge w\br{S}-\rho\cdot w\br{G},
\]
so every constraint is satisfied. The LP cost equals $c(F)$, thus
$\tau_X\le c(F)$ for every feasible $(\rho,X)$-separator $F$. Consequently,
$\tau_X$ is a lower bound on the optimum separator capacity.
\end{proof}

\subsection{Radius bound}

For a vertex $v\in S\subseteq V$, define $
\operatorname{radius}(v,S)=\max_{u\in S}\dist_S(v,u).$
We have the following lemma:

\begin{lemma}
If $S\subseteq V$ satisfies $w\br{S}>\rho_0\cdot w\br{G}$, then for every terminal $v\in S\cap X$, $
\operatorname{radius}(v,S)>\frac{\rho_0-\rho}{\rho_0}.
$
\end{lemma}

\begin{proof}
Fix $S$ and terminal $v\in S\cap X$. By \eqref{eq:px1-main}, we have $
\sum_{u\in S}\dist_{x,V}(v,u)\cdot w(u)\ge w\br{S}-\rho\cdot w\br{G}.
$
Hence the weighted average distance from $v$ to vertices in $S$ is at least
\[
1-\frac{\rho\cdot w\br{G}}{w\br{S}} > 1-\frac{\rho}{\rho_0}=\frac{\rho_0-\rho}{\rho_0}.
\]
By averaging there exists a vetrtex $u\in S$ such that $
\dist_{x,V}(v,u)>\frac{\rho_0-\rho}{\rho_0}.
$
Since $\dist_{x,S}(v,u)\ge \dist_{x,V}(v,u)$, we obtain
\[
\operatorname{radius}(v,S)\ge \dist_S(v,u)>\frac{\rho_0-\rho}{\rho_0}.
\]
\end{proof}

\subsection{Assigning volumes to spheres}

Fix a connected subgraph $G'=(V',E')$ that we currently process.
For a vertex $v\in V'$ and radius $r\ge 0$, define an \emph{$r$-sphere} as $
N(v,r)=\{u\in V'\colon \dist_{V'}(u,v)<r\}.$. Let $
E(v,r)=E'\cap \bigl(N(v,r)\times N(v,r)\bigr) $ be the set of edges with both endpoints in the $r$-sphere, and let $
\delta(v,r)=E'\cap \bigl(N(v,r),V'\setminus N(v,r)\bigr)$ be the set of edges with exactly one endpoint in the $r$-sphere.

For a parameter $\varepsilon>0$, define the volume of a sphere by
\[
\operatorname{vol}(v,r)=\varepsilon\tau_X
+\sum_{e\in E(v,r)} c(e)\cdot d(e)
+\sum_{(x,y)\in \delta(v,r)} c(x,y)\bigl(r-\dist_{N(v,r)}(v,x)\bigr).
\]

This quantity consists of 3 terms: a small constant term $\varepsilon\tau_X$ called \emph{seed value}, the contribution of edges inside the sphere, and the contribution of edges crossing the sphere boundary. Specifically, we will choose $\epsilon=\frac{1}{n\ln n}$.
Now, we observe that $\operatorname{vol}(v,r)$ is monotone and piecewise linear. This is because the first two terms are monotone and piecewise constant, and the last term is monotone and piecewise linear. In particular, $\operatorname{vol}(v,r)$ is continuous, and on each open interval between consecutive vertex distances from $v$,
\[
\frac{d}{dr}\operatorname{vol}(v,r)=c\bigl(\delta(v,r)\bigr).
\]

\subsection{Cut procedure and its correctness}

Here we describe the modified region-growing procedure, which is almost identical to the original procedure of Even et al.~\cite{FastApproximateGraphPartitioningAlgorithms}, except that it grows a region around an arbitrary terminal $v\in V'\cap X$ instead of an arbitrary vertex $v\in V'$. The procedure is given in Algorithm~\ref{alg:terminal-cut-proc}. This procedure is used, whenever there exists at least one terminal in the current subgraph $G'=(V',E')$ and $w\br{V'}>\rho_0\cdot w\br{G}$. The procedure returns a sphere $U=N(v,r)$ for some radius $r\le \tilde r$ such that $w\br{U}\le \rho_0\cdot w\br{G}$ and the ratio between $c\br{\delta\br{v, r}}$ and $\operatorname{vol}(v,r)$ is $O\br{\log n}$.

Pick $
\tilde r=\frac{\rho_0-\rho}{\rho_0}.$



\begin{algorithm}[H]
\caption{\textsc{Cut-Proc} for Terminal-Constrained Region Growing}
\label{alg:terminal-cut-proc}
\DontPrintSemicolon
\KwIn{A connected subgraph $G'=(V',E')$, edge capacities $c$, LP lengths $d$, terminal set $X\subseteq V'$, with $V'\cap X\neq\emptyset$ and $w(V')>\rho_0\cdot w(V)$}
\KwOut{A nontrivial set $U\subsetneq V'$ satisfying \eqref{eq:terminal-good-radius}}
Set $\tilde r\gets (\rho_0-\rho)/\rho_0$\;
Choose an arbitrary center $v\in V'\cap X$\;
Set $U\gets\{v\}$\;
Set $v_0\gets$ closest vertex to $U$ in $V'\setminus U$ with respect to $\dist_{d,V'}$\;
\While{$c\br{U, V'\setminus U} > \frac{1}{\tilde r}\cdot\ln\left(\frac{\operatorname{vol}(v,\tilde r)}{\operatorname{vol}(v,0)}\right)\cdot \operatorname{vol}\br{v,\dist_{d,V'}\br{v,v_0}}$}{
	$U\gets U\cup\{v_0\}$\;
	Set $v_0\gets$ closest vertex to $U$ in $V'\setminus U$ with respect to $\dist_{d,V'}$\;
}
\Return{$U$}\;
\end{algorithm}


\begin{theorem}
For every processed subgraph $G'=(V',E')$ with $V'\cap X\neq\emptyset$ and
$w\br{V'}>\rho_0\cdot w\br{G}$, procedure \textsc{Cut-Proc} returns
nonempty $U\subsetneq V'$ such that:
\begin{enumerate}
    \item $U\cap X\neq\emptyset$,
    \item $w\br{U}\le \rho_0\cdot w\br{G}$,
    \item The following inequality holds for the returned radius:
    \[
c\bigl(\delta(v,r)\bigr)
\le
\frac{1}{\tilde r}\ln\!\left(\frac{\operatorname{vol}(v,\tilde r)}{\operatorname{vol}(v,0)}\right)
\cdot \operatorname{vol}(v,r).
\label{eq:terminal-good-radius}
\]
\end{enumerate}
\end{theorem}

\begin{proof}
Fix the chosen terminal center $v\in V'\cap X$. Let
$z_0=v,z_1,\dots,z_{|V'|-1}$ be vertices sorted by nondecreasing
distance from $v$ in $G'$. For each $i$, denote
\[
I_i=\bigl(\dist_{V'}(v,z_i),\dist_{V'}(v,z_{i+1})\bigr).
\]
On each $I_i$, $\operatorname{vol}(v,r)$ is linear and
\[
\operatorname{vol}'(v,r)=c\bigl(\delta(v,r)\bigr)
\qquad\text{for }r\in I_i.
\]

\begin{lemma}\label{lem:cut-proc-log-derivative-point}
There exists $r_0\in (0,\tilde r]\cap\bigcup_i I_i$ such that
\[
\frac{\operatorname{vol}'(v,r_0)}{\operatorname{vol}(v,r_0)}
\le
\frac{1}{\tilde r}
\ln\!\left(\frac{\operatorname{vol}(v,\tilde r)}{\operatorname{vol}(v,0)}\right).
\]
\end{lemma}
\begin{proof}
We use an averaging argument for the logarithmic derivative over $[0,\tilde r]$.
Assume to the contrary that for all admissible $r\in(0,\tilde r]$,
\[
\frac{\operatorname{vol}'(v,r)}{\operatorname{vol}(v,r)}
>
\frac{1}{\tilde r}
\ln\!\left(\frac{\operatorname{vol}(v,\tilde r)}{\operatorname{vol}(v,0)}\right).
\]
Integrating over $[0,\tilde r]$ gives
\[
\int_0^{\tilde r}\frac{\operatorname{vol}'(v,r)}{\operatorname{vol}(v,r)}dr
>
\frac{1}{\tilde r}
\ln\!\left(\frac{\operatorname{vol}(v,\tilde r)}{\operatorname{vol}(v,0)}\right)
\int_0^{\tilde r}dr.
\]
Both sides equal
\[
\ln\operatorname{vol}(v,\tilde r)-\ln\operatorname{vol}(v,0)
=
\ln\!\left(\frac{\operatorname{vol}(v,\tilde r)}{\operatorname{vol}(v,0)}\right),
\]
contradiction.
\end{proof}

\begin{lemma}\label{lem:cut-proc-scanned-radius}
There exists a scanned radius $r\le \tilde r$ satisfying \eqref{eq:terminal-good-radius}.
\end{lemma}
\begin{proof}
We first find a good continuous radius and then move to the nearest scanned event radius, keeping the same cut while not increasing the relevant ratio.
Let $r_0$ be given by Lemma~\ref{lem:cut-proc-log-derivative-point}, and define
\[
r_1=\min\{\dist_{V'}(v,u)\colon u\in V',\ \dist_{V'}(v,u)\ge r_0\}.
\]
Then $N(v,r_0)=N(v,r_1)$, hence both radii induce the same cut and
\[
c\bigl(\delta(v,r_1)\bigr)=\operatorname{vol}'(v,r_0).
\]
By monotonicity $\operatorname{vol}(v,r_0)\le \operatorname{vol}(v,r_1)$, so \eqref{eq:terminal-good-radius} holds for $r_1$. Since the procedure scans event radii exhaustively, it returns some scanned radius $r\le r_1\le \tilde r$ satisfying \eqref{eq:terminal-good-radius}.
\end{proof}

\begin{lemma}\label{lem:cut-proc-weight-bound}
For the returned set $U=N(v,r)$, we have $w\br{U}\le \rho_0\cdot w\br{G}$.
\end{lemma}
\begin{proof}
We prove this by contradiction: if the returned ball were too heavy, the radius lower bound for terminal-centered sets would force a vertex farther than $\tilde r$ from $v$ inside $U$.
If $w\br{U}>\rho_0\cdot w\br{G}$, then by terminal radius guarantee applied to $S=U$ and center $v\in U\cap X$,
\[
\operatorname{radius}(v,U)>\tilde r.
\]
But $U=N(v,r)$ and $r\le \tilde r$, so every $u\in U$ satisfies $\dist_{V'}(v,u)<r\le \tilde r$, contradiction.
Hence $w\br{U}\le \rho_0\cdot w\br{G}$.
\end{proof}

By Lemma~\ref{lem:cut-proc-scanned-radius}, property (3) of the theorem holds. By Lemma~\ref{lem:cut-proc-weight-bound}, property (2) holds. Property (1) is immediate since $v\in U\cap X$.
\end{proof}

\subsection{Global algorithm and approximation factor}

Run \textsc{Cut-Proc} iteratively as in Algorithm~\ref{alg:terminal-global-procedure}.

\begin{algorithm}[H]
\caption{Global Terminal Separator via Iterated \textsc{Cut-Proc}}
\label{alg:terminal-global-procedure}
\DontPrintSemicolon
\KwIn{A weighted graph $G=(V,E,c,w)$, terminal set $X\subseteq V$, and parameters $0<\rho<\rho_0<1$}
\KwOut{A feasible $(\rho_0,X)$-separator $F\subseteq E$}
Initialize $F\gets\emptyset$\;
\While{there exists a connected component $C$ of $(V,E\setminus F)$ with $C\cap X\neq\emptyset$ and $w(C)>\rho_0 w(V)$}{
    Run Algorithm~\ref{alg:terminal-cut-proc} on $G[C]$ with terminal set $X\cap C$ and obtain $U\subsetneq C$\;
    Add boundary edges of $U$ in $G[C]$ to $F$, i.e.,
    $F\gets F\cup \delta_{G[C]}(U).$
}
\Return{$F$}\;
\end{algorithm}


At termination, every terminal-containing component has weight at most $\rho_0\cdot w\br{G}$, so the produced cut is a feasible $(\rho_0,X)$-separator. 
Let the procedure be called $q$ times and let $U_1,\dots,U_q$ be returned spheres.
By construction, these sets are pairwise disjoint.

For each call, by \eqref{eq:terminal-good-radius},
\[
c\bigl(\delta(U_j)\bigr)
\le
\frac{1}{\tilde r}
\ln\!\left(\frac{\operatorname{vol}_j(v_j,\tilde r)}{\operatorname{vol}_j(v_j,0)}\right)
\operatorname{vol}_j(v_j,r_j),
\]
where $\operatorname{vol}_j$ is the volume function within the current processed component.
Since
\[
\operatorname{vol}_j(v_j,\tilde r)\le (1+\varepsilon)\tau_X,
\qquad
\operatorname{vol}_j(v_j,0)=\varepsilon\tau_X,
\]
we get
\[
c\bigl(\delta(U_j)\bigr)
\le
\frac{1}{\tilde r}\ln\!\left(\frac{1+\varepsilon}{\varepsilon}\right)
\operatorname{vol}_j(v_j,r_j).
\]

Summing over calls, edge-length contributions are at most $\tau_X$, and seed terms contribute at most $q\varepsilon\tau_X\le n\varepsilon\tau_X$. Hence
\[
\sum_{j=1}^q c\bigl(\delta(U_j)\bigr)
\le
\frac{1}{\tilde r}
\ln\!\left(\frac{1+\varepsilon}{\varepsilon}\right)
\tau_X(1+\varepsilon n).
\]

Choosing
\[
\varepsilon=\frac{1}{n\ln n},
\]
we obtain the following bound.

\begin{theorem}
The above algorithm computes a $(\rho_0,X)$-separator of capacity at most
\[
\left(\frac{\rho_0}{\rho_0-\rho}+o(1)\right)\ln n\cdot \tau_X.
\]
\end{theorem}
